\subsection{Curve contributions and global places}
\label{subsec:curve-places}

It remains to bound the contribution \(b_T\) of each singular curve by
global divisors.  We realize the relevant cohomology classes by line
bundles near a compact subset carrying all loops of \(T^\circ\).  Degrees
on the fibers of a small \(\Q\)-factorialization then produce exceptional
surfaces.  Retaining those loops will ensure that their blowup places
remain distinct on the fixed global resolution.

\begin{proposition}[Global places over a singular curve]\label{prop:curve-places}
For every irreducible curve component \(T\) of \(\Sing Y\), with \(b_T\)
as in \eqref{eq:branch-local-systems} and the fixed resolution
\(g:\widetilde Y\to Y\),
\[
 b_T\ \le\
 \#\left\{
 \begin{array}{c}
 E\text{ a prime component of }\Exc(g):\\
 g(E)=T,\quad a(E;Y,0)=2
 \end{array}
 \right\}.
\]
In particular, the divisors counted here represent distinct global
divisorial places.
\end{proposition}

\begin{proof}
Fix \(T\), and write
\[
 q_T:\mathcal W_T\longrightarrow\mathcal V_T
\]
for the natural quotient of Lemma~\ref{lem:link-quotient}.  If \(b_T=0\)
there is nothing to prove, so assume \(b_T>0\).  We may remove finitely
many further points of \(T^\circ\) whenever needed.  The local systems
already defined restrict to the smaller curve, and the induced map on
fundamental groups is surjective.  Thus their monodromy images and the
coinvariant dimension \(b_T\) do not change.

We first construct sections
\[
 \sigma_1,\ldots,\sigma_{b_T}\in H^0(T^\circ,\mathcal W_T)
\]
whose images under \(q_T\) are linearly independent in every fiber of
\(\mathcal V_T\), after the permitted further finite puncturing.

\emph{Step 1. Invariant lifts from finite monodromy.}
First we show that \(\mathcal W_T\), and therefore its quotient
\(\mathcal V_T\), has finite monodromy.  Let
\(\mathcal F=\widetilde Y\times_Y T\), with its full complex-space
structure.  It is projective over the compact projective curve \(T\).
A relative embedding and GAGA identify this morphism with an algebraic
projective family, including its possibly nonreduced scheme fibers.
After deleting finitely many points we may assume that
\(\mathcal F\to T^\circ\) is flat and that all
\((R^kg_*\Q)|_{T^\circ}\) are local systems.

For \(t\in T^\circ\), put \(F_t=\mathcal F_t\).  Every integral class in
\(H^2(F_t,\Z)\) lifts, by proper base change, to
\(H^2(g^{-1}U_t,\Z)\) for a sufficiently small Stein neighborhood
\(U_t\) of \(t\) in \(Y\).  Terminal singularities are rational
\cite[Remark~3.8 and Theorem~3.12]{FujinoAnalyticMMP}.  Hence
\(g_*\mathcal O_{\widetilde Y}=\mathcal O_Y\) and
\(R^kg_*\mathcal O_{\widetilde Y}=0\) for \(k>0\); Leray and Cartan's
Theorem~B give
\[
 H^2(g^{-1}U_t,\mathcal O_{\widetilde Y})=0.
\]
The exponential sequence therefore realizes the lifted integral class
as the first Chern class of a holomorphic line bundle on \(g^{-1}U_t\).
Its restriction to the projective scheme \(F_t\) is algebraic by the
scheme form of GAGA
\cite[Corollary~4.3 and Theorem~4.4]{RaynaudSGA1XII}.  In particular,
\begin{equation}\label{eq:curve-fiber-chern}
 \Pic(F_t)\otimes_{\Z}\Q
 \xrightarrow{\ c_1\ } H^2(F_t,\Q)=(\mathcal W_T)_t
 \quad\text{is surjective}.
\end{equation}
Here and below the cohomology of a scheme fiber means that of its
analytic topological space; nilpotents do not change it.

We use the parameter argument of
\cite[Section~6.1, proof of Lemma~6.4]{ProjectiveLCFourfolds} to show
that a spanning set of fiber classes has finite monodromy orbits.  Fix a
relatively very ample \(\mathcal O_{\mathcal F}(1)\).  Every line bundle on a
fiber is a quotient of some
\(\mathcal O_{F_t}(-n)^{\oplus r}\).  Letting \(n,r\), and the Hilbert
polynomial vary gives countably many relative Quot schemes of finite
type over \(T^\circ\)
\cite[Theorems~3.1--3.2]{GrothendieckHilbert1961}.  In each take the
open set on which the universal quotient is invertible on the whole
pulled-back family.  These open sets still contain a point representing
every fiber line bundle.  Indeed the family and universal quotient are
flat over the Quot scheme, so the local criterion for flatness makes a
fiberwise locally free quotient of rank one locally free near that
fiber; properness removes the image of the complementary closed locus.

There are countably many irreducible components of these parameter
spaces.  Choose \(t\) outside the closures of the images of all
components which do not dominate \(T^\circ\).  Such a point exists,
since each of those closures is a finite set and a complex curve is
uncountable.  Put \(\Gamma=\pi_1(T^\circ,t)\).  By
\eqref{eq:curve-fiber-chern}, choose finitely many
line bundles on \(F_t\) whose Chern classes span
\((\mathcal W_T)_t\).  Each is represented on a dominating
reduced irreducible parameter component \(H\).  The Chern class of the
universal bundle gives a section of the pullback of \(\mathcal W_T\)
to \(H^{\mathrm{an}}\): restrict its class on the total family to
fibers and use proper base change.

Choose a closed point of the generic fiber of \(H\to T^\circ\).
Its residue field is a finite extension of the function field of
\(T\).  Spreading this point, normalizing its closure, and deleting
finitely many base points gives a connected finite \'etale cover of a
dense open of \(T^\circ\), together with a morphism from that cover
to \(H\).  The analytification of an
irreducible complex algebraic variety is path connected.  A path in
\(H^{\mathrm{an}}\) from the original parameter point to a point of
this cover identifies the transported Chern class with the class
there.  The subgroup fixing that sheet has finite index and fixes
this class.  The inclusion of the further punctured curve, with the
projected path for change of basepoint, still surjects on \(\Gamma\).
The image of the finite-index sheet stabilizer is therefore finite
index in \(\Gamma\) and fixes the original class.  Thus each of our
spanning classes has a finite orbit under \(\Gamma\).  The intersection
of their stabilizers has finite index and acts trivially on
\((\mathcal W_T)_t\).  The monodromy group \(G\) of \(\mathcal W_T\)
is therefore finite, as is that of \(\mathcal V_T\).

Poincar\'e duality with local coefficients on the oriented curve gives
\[
 b_T=\dim(\mathcal V_T)_t{}_{\Gamma}.
\]
For a representation of the finite group \(G\) over \(\Q\), averaging
identifies coinvariants with invariants and makes invariants an exact
functor.  The quotient \(q_T\) therefore supplies the asserted sections
\(\sigma_i\) with independent images.

\emph{Step 2. Line bundles near a compact curve core.}
Take \(T^\circ\) noncompact, removing one more point if needed.
Recall the fixed choice made after \eqref{eq:branch-local-systems}:
for every prime component \(E\) of \(\Exc(g)\) dominating \(T\),
its Stein factor restricts to
\begin{equation}\label{eq:global-divisor-family}
 E|_{T^\circ}\longrightarrow \widehat T_E^\circ\longrightarrow T^\circ,
\end{equation}
where the first map is a smooth proper family with connected fibers
and the second is a connected finite \'etale cover.  Further finite
puncturing preserves these properties.  Choose an ambient open
\(Y^\circ\subset Y\) in which \(T^\circ\) is closed: remove the
discarded points of \(T\), the other singular components, and their
intersection points.  These choices will be used when we identify the
local divisors globally.

Let \(h:g^{-1}T^\circ\to T^\circ\) be the restricted resolution.
The curve \(T^\circ\) has the homotopy type of a finite graph, so
cohomology of its local systems vanishes in degrees at least two.
The Leray spectral sequence for \(h\) consequently makes the edge map
\[
 H^2(g^{-1}T^\circ,\Q)
 \longrightarrow H^0(T^\circ,\mathcal W_T)
\]
surjective.  Choose lifts of the \(\sigma_i\).  The algebraic space
\(g^{-1}T^\circ\) has finite triangulation type; after multiplying
each lift by a positive integer, it comes from a class
\(\alpha_i\in H^2(g^{-1}T^\circ,\Z)\).  Replace \(\sigma_i\) by the
same multiple.  Their images are still independent.

To extend the \(\alpha_i\) simultaneously while retaining every loop
of \(T^\circ\), choose a compact curve core as follows.  In the smooth
compactification of \(T^\circ\), remove small open
coordinate discs about its punctures, with disjoint closures, and denote the complement by
\(K_*\).  This is a connected compact semianalytic subset of
\(T^\circ\).  The inclusion of its interior into \(T^\circ\) induces
a surjection on fundamental groups.  The core \(K_*\) is also Runge: every component of
\(T^\circ\setminus K_*\) runs to a missing puncture and is therefore
not relatively compact.  The compact set \(g^{-1}K_*\) is
semianalytic.  Continuity of cohomology on neighborhoods of this
compact set extends the restrictions of the finitely many
\(\alpha_i\) to integral classes on one open neighborhood
\(\widetilde U_0\) of \(g^{-1}K_*\) in \(\widetilde Y\).
One may use \v{C}ech cohomology for this continuity statement;
semianalytic triangulation and local contractibility identify it with
singular cohomology here.  Properness of \(g\) makes
\[
 U_0=Y\setminus g(\widetilde Y\setminus\widetilde U_0)
\]
an open neighborhood of \(K_*\) with \(g^{-1}U_0\subset\widetilde U_0\).

The closed curve \(T^\circ\subset Y^\circ\) is Stein.  Siu's
Stein-neighborhood theorem, in the form
\cite[Theorem~2.1]{ForstnericStein} with empty compact set, gives a
Stein neighborhood \(\Omega\) of this curve in \(Y^\circ\).
The Runge property makes \(K_*\) holomorphically convex in
\(T^\circ\), hence also in \(\Omega\) by the paragraph preceding
that theorem.  The usual holomorphic-polyhedron construction for a
convex compact in a Stein space gives a Stein neighborhood
\[
 K_*\subset U\Subset\Omega\cap U_0.
\]
Take its connected component containing \(K_*\).  It is normal,
being open in \(Y\), and therefore irreducible.  For the local
\(\Q\)-factorialization, Fujino's compact construction
\cite[Section~1.11]{FujinoAnalyticMMP} supplies an ambient semianalytic
Stein compact \(K'\Subset U\), with
\(K_*\subset\operatorname{int}_U(K')\), such that
\[
 K'\cap Z\ \text{has finitely many connected components}
\]
for every analytic subset \(Z\) defined on a neighborhood of
\(K'\).  This is the compact condition (P4) in Fujino's condition
(P); normality of the source and Steinness of the base supply its
other ambient conditions.

The classes just extended to \(g^{-1}U\) are Chern classes of line
bundles \(L_1,\ldots,L_{b_T}\) there, by the same rational-singularity,
Cartan, and exponential-sequence argument used above.  For every
\(t\) in the interior of \(K_*\), their fiber classes are the
\(\sigma_i(t)\); hence their images under \(q_T\) remain independent.
We also need divisor representatives before making any further
modification.  The sheaf \(g_*L_i\) is coherent and has rank one on
the regular open of \(U\).  Cartan's Theorem~A supplies a global
section nonzero at a point of that open.  Its pullback is a nonzero
section of \(L_i\), whose zero divisor is a Cartier divisor \(A_i\)
on \(g^{-1}U\) representing \(L_i\).

\emph{Step 3. Fiber degrees and surface families.}
We will obtain at least \(b_T\) surface families by proving that the
\(b_T\) line bundles constructed below have linearly independent vectors
of fiber degrees.
Apply Fujino's small \(\Q\)-factorialization theorem
\cite[Theorem~1.24]{FujinoAnalyticMMP} to the ordinary pair
\((U,0)\), the identity \(U\to U\), and the full compact \(K'\).
The identity is projective, \(U\) is terminal and therefore klt,
and \(K_U\) is \(\Q\)-Cartier.  The preceding construction supplies
condition (P).  After restricting the base to a neighborhood of
\(K'\), the theorem gives a projective small bimeromorphic morphism
\[
 \ell:Y'\longrightarrow U
\]
with \(Y'\) normal and \(\Q\)-factorial over \(K'\).  We retain the
notation \(U\) after this restriction and restrict \(g,L_i,A_i\)
accordingly.  Compactness and local terminality of \(Y\) give a common
\(m_{\mathrm{can}}>0\) for which \(\omega_Y^{[m_{\mathrm{can}}]}\) is
invertible; its restriction to \(U\) is
\(\omega_U^{[m_{\mathrm{can}}]}\).  Since \(\ell\) is
small, the natural identification on the common codimension-one open
extends by reflexivity to
\[
 \omega_{Y'}^{[m_{\mathrm{can}}]}
 \simeq\ell^*\omega_U^{[m_{\mathrm{can}}]}.
\]
In particular, the left side is invertible and the natural relative
canonical divisor \(K_{Y'/U}\) is zero.  This is canonical crepancy,
expressed using compatible local canonical generators as
\(K_{Y'}=\ell^*K_U\).  A place exceptional over \(Y'\) is exceptional
over \(U\), because \(\ell\) is small.  Crepancy therefore preserves
the terminal discrepancy inequalities, so \(Y'\) is terminal.  Moreover \(\ell\)
is an isomorphism over \(U_{\mathrm{reg}}\).  Indeed, for a relatively
ample line bundle \(\mathcal A\) on \(Y'\), the rank-one reflexive
sheaf \((\ell_*\mathcal A)^{**}\) is locally invertible on the smooth
base.  Its pullback agrees with \(\mathcal A\) off codimension two
and hence everywhere there, by reflexivity on the normal source.
A positive-dimensional fiber over a smooth point would then have
degree zero against \(\mathcal A\), contradicting relative ampleness.

Push \(A_i\) forward to a Weil divisor on \(U\), and let \(D_i'\) be
its strict transform on \(Y'\).  These are actual locally finite
analytic divisors defined on all of the current \(Y'\).  We apply
\(\Q\)-factoriality precisely in the sense of
\cite[Definition~2.38]{FujinoAnalyticMMP}: every prime divisor
defined on a neighborhood of the whole set \(\ell^{-1}(K')\)
is \(\Q\)-Cartier at each point of that set.  Properness makes
\(\ell^{-1}(K')\) compact, and local finiteness leaves only finitely
many prime components of the \(D_i'\) near it.  A finite cover of
this full compact and a common multiple therefore give an integer
\(m_{\mathrm{div}}>0\) and an open neighborhood of \(\ell^{-1}(K')\)
on which every \(m_{\mathrm{div}}D_i'\) is Cartier.  By properness,
choose an open set \(U_1\) with \(K'\subset U_1\subset U\) whose full
inverse image is contained in that neighborhood, and put
\[
 \mathcal N_i=\mathcal O_{Y'}(m_{\mathrm{div}}D_i')|_{\ell^{-1}U_1}.
\]
This completes the use of \(\Q\)-factoriality on the specified full
compact; all later restrictions use these fixed line bundles.
In particular, the change-of-compact issue in
\cite[Remark~2.39]{FujinoAnalyticMMP} does not arise.  On the
regular open, where both \(g\) and \(\ell\) are isomorphisms,
\begin{equation}\label{eq:curve-regular-bundles}
 \mathcal N_i|_{U_{1,\mathrm{reg}}}
 \simeq L_i^{\otimes m_{\mathrm{div}}}|_{U_{1,\mathrm{reg}}}.
\end{equation}

Choose connected curve domains
\[
 B_0\Subset B_1\Subset\operatorname{int}_{T^\circ}(K_*)
\]
such that \(\pi_1(B_0)\to\pi_1(T^\circ)\) is surjective; nested
smaller cores obtained by enlarging the omitted discs have this
property.  The larger domain lets us control finitely many components
near \(\overline B_0\) and make their discrete exceptional sets finite
on that compact.  Work in the ambient open of \(U_1\) obtained by removing
\(T^\circ\setminus B_1\), so that \(B_1\) is closed there.
The inverse image of \(B_1\) is proper over \(B_1\).
Local finiteness of its analytic components, compactness of
\(\ell^{-1}(\overline B_0)\), and properness allow us first to
replace \(B_1\) by a neighborhood of \(\overline B_0\) on whose
inverse image only finitely many components occur.

We claim that deleting finitely many points from \(B_0\) gives a
connected curve domain \(B\) with
\(\pi_1(B)\to\pi_1(T^\circ)\) surjective and the following fiber
description.  All fibers of \(\ell\) over \(B\) have dimension at
most one.  The finitely many irreducible surface components
\(S_1,\ldots,S_s\) of the reduced inverse image of \(B\), counted
after this restriction, all dominate \(B\).  Their normalizations
factor as
\[
 \widehat S_\alpha\xrightarrow{\,h_\alpha\,}C_\alpha
 \longrightarrow B \qquad (1\le\alpha\le s),
\]
where \(C_\alpha\to B\) is a connected finite \'etale cover and
\(h_\alpha\) is a smooth proper family of connected curves.  A cover
that splits on restriction contributes separate components to this
list.  For each fixed \(t\in B\), the pairs \((\alpha,c)\) with
\(c\in C_\alpha\) over \(t\) supply distinct irreducible curves and
exhaust the irreducible curves of \(\ell^{-1}(t)\).  Each smooth
connected normalization fiber maps birationally to the curve it
supplies.

We verify this description by spelling out the finite exclusions.
Smallness gives \(\dim\Exc(\ell)\le2\), so a general fiber above
\(B_1\) has dimension at most one.  The locus of larger fibers is a
proper analytic subset of \(B_1\), by the fiber-dimension theorem
for proper maps.  For each surface component of the reduced inverse
image that dominates \(B_1\), take its normalization and then its Stein factor over the
curve.  The normalization is proper and is a normal surface.
Its singularities are isolated.  Away from the branch values of the
finite Stein factor and the critical values of the map from the
normalization, it is a smooth proper family of connected curves
over a finite unramified cover.  The values at which an intersection
of two components, or the nonisomorphism locus of a normalization,
contains a whole fiber curve also form a discrete analytic subset:
these loci have dimension at most one, and their general fibers
over the base are zero-dimensional.  Finally remove the images of
components lying over a point.  Each of these excluded sets is
discrete on the larger curve domain \(B_1\); there are only finitely
many components, so their union meets \(\overline B_0\) in a finite
set.

Delete those finitely many points from \(B_0\) and call the result
\(B\).  It is connected, and
\(\pi_1(B)\to\pi_1(T^\circ)\) is still surjective: loops in \(B_0\)
may be perturbed away from a finite set.  The exclusions give the
asserted fiber description and families.

For each \(\alpha,i\), let
\[
 d_{\alpha i}
 =\deg\bigl(\mathcal N_i|_{h_\alpha^{-1}(c)}\bigr),
 \qquad c\in C_\alpha,
\]
where pullback to \(\widehat S_\alpha\) is understood.  This integer
is independent of \(c\), by constancy of degree in a smooth proper
family over the connected curve \(C_\alpha\).  Because its fiber
maps birationally to the curve it supplies, this is also the degree
of \(\mathcal N_i\) on that fiber curve of \(\ell\).
We claim that the map
\begin{equation}\label{eq:curve-degree-map}
 \Q^{b_T}\longrightarrow\Q^s,\qquad
 (c_i)_i\longmapsto
 \left(\sum_i c_i d_{\alpha i}\right)_\alpha
\end{equation}
is injective.

Suppose its value is zero, and fix \(t\in B\).  Every irreducible
curve in \(F'_t=\ell^{-1}(t)\) then has degree zero against
\(\sum_i c_i c_1(\mathcal N_i)\).  The fiber has dimension at most
one.  Its rational \(H_2\) is spanned by the fundamental classes of
its irreducible curves, also when the fiber is nonreduced, and
\(H^2(F'_t,\Q)\) is the rational dual of \(H_2(F'_t,\Q)\).
The restricted Chern class is therefore zero.  Proper base change
identifies it with the germ in \((R^2\ell_*\Q)_t\).
Restriction through the isomorphism over the regular open gives
the natural map
\[
 (R^2\ell_*\Q)_t\longrightarrow (R^2j_*\Q)_t
     =(\mathcal V_T)_t,
\]
so the corresponding punctured class is zero as well.  By
\eqref{eq:curve-regular-bundles}, the restrictions of
\(\sum_i c_i c_1(\mathcal N_i)\) and
\(m_{\mathrm{div}}\sum_i c_i c_1(L_i)\) agree on their common regular
open.  Taking their germ in \((R^2j_*\Q)_t\) and applying
\eqref{eq:branch-link-stalk-map} to the latter identifies that zero
class with
\[
 m_{\mathrm{div}}\sum_i c_i\,q_T(\sigma_i(t)).
\]
These vectors are independent; hence every \(c_i=0\).  This proves
\eqref{eq:curve-degree-map}, and in particular
\begin{equation}\label{eq:curve-surface-count}
 s\ge b_T.
\end{equation}

\emph{Step 4. Distinct global discrepancy-two divisors.}
Take an ambient open \(U_B\subset U_1\) whose intersection with
\(T^\circ\) is the closed subspace \(B\), obtained as above by
removing \(T^\circ\setminus B\).  Each \(S_\alpha\) is a closed
analytic surface in \(\ell^{-1}U_B\).  Since \(Y'\) is terminal,
it is smooth in codimension two by
\cite[Lemma~2.3]{CanonicalKahler}.  Thus at a general point of
\(S_\alpha\) both the ambient space and the surface are smooth.
Blow up its coherent ideal and normalize.  The unique divisor
dominating that dense smooth part defines a local place \(v_\alpha\)
of ordinary log discrepancy \(2\) over \(Y'\), the value for a
smooth codimension-two blowup.  Canonical crepancy of \(\ell\)
gives
\[
 a(v_\alpha;U_B,0)=2,\qquad
 \cent_{U_B}(v_\alpha)=B.
\]
The \(v_\alpha\) are distinct because their centers on \(Y'\) are
the distinct \(S_\alpha\).

It remains to represent each \(v_\alpha\) by a prime divisor on
\(g^{-1}U_B\) and show that these local divisors arise from distinct
global components of \(\Exc(g)\).

On the fixed global resolution, the relative canonical divisor is
intrinsic even if the canonical bundle has no global divisor
representative.  Using the fixed \(m_{\mathrm{can}}\), divide by
\(m_{\mathrm{can}}\) the divisor of the
natural meromorphic comparison
\(g^*\omega_Y^{[m_{\mathrm{can}}]}\dashrightarrow
\omega_{\widetilde Y}^{\otimes m_{\mathrm{can}}}\), and denote the
result by \(K_{\widetilde Y/Y}\).  It is exceptional; write
\[
 K_{\widetilde Y/Y}=\sum_E\kappa_EE.
\]
The natural canonical comparisons commute with restriction and
composition, so this divisor supplies the relative discrepancy term
on \(U_B\) as well.
Every \(\kappa_E\) is strictly positive by terminality.  The
principalization of the ideal of \(\Sing Y\) says that
\(g^{-1}(\Sing Y)\) is supported on these exceptional divisors.
Put \(\widetilde U_B=g^{-1}U_B\), and consider \(v_\alpha\) on a
common proper model with \(\widetilde U_B\).  If it were exceptional
over this smooth resolution, then \(a(v_\alpha;\widetilde U_B,0)\ge2\).
Its center on the
resolution maps into \(B\subset\Sing Y\), so it is contained in at
least one of the \(E\), with \(\ord_{v_\alpha}(E)>0\).
The discrepancy formula would give
\[
 a(v_\alpha;U_B,0)
 =a(v_\alpha;\widetilde U_B,0)
    +\sum_E\kappa_E\ord_{v_\alpha}(E)>2,
\]
a contradiction.  Thus \(v_\alpha\) is represented by a prime
divisor on \(g^{-1}U_B\).  Since \(U_B\cap T=B\), this local divisor
is an irreducible component of \(E|_B\) for a global exceptional
divisor \(E\) with \(g(E)=T\).

Such a global divisor cannot split into two components over \(B\).
Indeed the connected finite cover \(\widehat T_E^\circ\to T^\circ\) in
\eqref{eq:global-divisor-family} remains connected over \(B\):
the action of \(\pi_1(T^\circ)\) on a fiber is transitive, and
\(\pi_1(B)\) has the same image.  Over this connected cover the
restriction of \(E\) is a smooth proper family with connected
fibers, so its total space is connected and smooth, hence
irreducible.  Therefore the distinct local divisors representing
the \(v_\alpha\) come from distinct global components \(E\) of
\(\Exc(g)\).  Their global centers are \(T\), and their ordinary
log discrepancies are \(2\), since this equality holds on the
open over \(B\).  Together with \eqref{eq:curve-surface-count},
this proves the proposition.
\end{proof}

\begin{proof}[Proof of Theorem~\ref{thm:branch}]
Every global place of ordinary log discrepancy \(2\) whose center
is a singular curve is represented by a component of \(\Exc(g)\).
Indeed the same discrepancy computation in the last part of
Proposition~\ref{prop:curve-places} would otherwise make its
discrepancy strictly greater than \(2\).  There are only finitely
many such components on the compact resolution, so
\(\tau_Y(Q)\) is finite.  Proposition~\ref{prop:curve-places},
summed over the singular curves contained in \(Q\), gives
\[
 \sum_{T\subset\Sing Y\cap Q} b_T\le\tau_Y(Q).
\]
Combining this with \eqref{eq:branch-topological-bound} yields
\[
 \beta(Q)\le \sum_j c_2(Q_j^\nu)+\tau_Y(Q),
\]
as required.
\end{proof}
